% ═══════════════════════════════════════════════════════════════════
%  Rough Volatility Arbitrage under Markov Regime Switching
%  A Volterra Process Approach
%  IEEEtran conference format — full manuscript
% ═══════════════════════════════════════════════════════════════════
\documentclass[conference]{IEEEtran}

\usepackage[utf8]{inputenc}
\usepackage[T1]{fontenc}
\usepackage[numbers]{natbib}


\usepackage{amsmath,amssymb,amsfonts,amsthm}
\usepackage{graphicx}
\usepackage{textcomp}
\usepackage{xcolor}
\usepackage{booktabs}
\usepackage{hyperref}
\usepackage{algorithm}
\usepackage{algpseudocode}
\usepackage{microtype}
\usepackage{enumitem}
\usepackage{siunitx}
\usepackage{multirow}
\usepackage{listings}
\usepackage{bm}
\usepackage{mathtools}

\hypersetup{colorlinks=true,linkcolor=black,urlcolor=blue,citecolor=black}

% ── Theorem environments ────────────────────────────────────────────
\newtheorem{remark}{Remark}
\newtheorem{proposition}{Proposition}
\newtheorem{definition}{Definition}

% ── Notation shortcuts ──────────────────────────────────────────────
\newcommand{\E}{\mathbb{E}}
\newcommand{\R}{\mathbb{R}}
\newcommand{\Prob}{\mathbb{P}}
\newcommand{\F}{\mathcal{F}}
\newcommand{\N}{\mathcal{N}}
\newcommand{\G}{\Gamma}
\newcommand{\dd}{\mathrm{d}}
\newcommand{\dW}{\,\dd W}
\newcommand{\abs}[1]{\lvert #1 \rvert}
\DeclareMathOperator{\Var}{Var}
\DeclareMathOperator{\Cov}{Cov}

% ── Code listing style ──────────────────────────────────────────────
\definecolor{codegreen}{rgb}{0,0.55,0}
\definecolor{codegray}{rgb}{0.45,0.45,0.45}
\definecolor{codepurple}{rgb}{0.58,0,0.82}
\definecolor{backcolour}{rgb}{0.97,0.97,0.97}
\lstdefinestyle{pystyle}{
    backgroundcolor=\color{backcolour},
    commentstyle=\color{codegreen},
    keywordstyle=\color{blue},
    numberstyle=\tiny\color{codegray},
    stringstyle=\color{codepurple},
    basicstyle=\ttfamily\footnotesize,
    breaklines=true,
    captionpos=b,
    keepspaces=true,
    numbers=left,
    numbersep=5pt,
    language=Python,
    morekeywords={njit,prange}
}
\lstset{style=pystyle}

% ═══════════════════════════════════════════════════════════════════
\begin{document}

\title{Rough Volatility Arbitrage under Markov Regime\\
       Volterra Process Approach with Double Exponential}

\author{
  \IEEEauthorblockN{Mitchell Scott, Ph.D.}
  \IEEEauthorblockA{
    \textit{Department of Mathematics}\\
    Emory University\\
    \text{mtscot4@emory.edu}
  }
  \and
  \IEEEauthorblockN{Felipe Cardozo}
  \IEEEauthorblockA{
    \textit{Mathematics \& Computer Science}\\
    Emory University\\
    \text{focardo@emory.edu}
  }
}

\maketitle

% ═══════════════════════════════════════════════════════════════════
\begin{abstract}
It is well-established in the empirical literature that the realised
variance of equity indices exhibits \emph{rough} behaviour characterised
by a Hurst exponent $H \approx 0.07$, substantially lower than the
value of $0.5$ implied by classical diffusion models.  The present
work introduces a unified trading framework—termed
\textit{Project Rough-Regime}—that (i) prices short-dated volatility
derivatives using a discrete Volterra stochastic process whose spot
variance is driven by a rough fractional kernel, (ii) classifies the
prevailing market regime in real time via a two-state Gaussian Hidden
Markov Model (G-HMM) calibrated by the Expectation–Maximisation
algorithm, and (iii) routes executions through an Interactive Brokers
TWS interface equipped with a passive-execution \emph{chase} algorithm.
The combined system constitutes an event-driven ``cyborg'' loop in
which the G-HMM serves as a \emph{traffic light} that suppresses long-volatility
exposure whenever the posterior probability of the turbulent state
exceeds $0.6$.  Empirical results obtained from Monte Carlo simulation
suggest that the regime-adjusted strategy achieves a Sharpe ratio
superior to a naive long-volatility benchmark by a factor of
approximately $1.4\times$, while the maximum drawdown is reduced by
more than $40\%$ at representative transaction-cost levels.  The
theoretical convergence of the Monte Carlo estimator at the expected
rate of $O(N^{-1/2})$ is verified experimentally, and the stability
of the HMM state-detection mechanism is validated against synthetic
market crises.
\end{abstract}

\begin{IEEEkeywords}
rough volatility, fractional Brownian motion, Volterra process, hidden
Markov model, Baum--Welch, Viterbi, Kelly criterion, high-frequency
trading, Interactive Brokers, Monte Carlo convergence, Heston model,
SABR model, local volatility, Dupire equation, Hawkes process,
jump-diffusion, stochastic volatility taxonomy.
\end{IEEEkeywords}

% ═══════════════════════════════════════════════════════════════════
\section{Introduction}\label{sec:intro}

The observation that equity-index volatility exhibits rougher sample
paths than classical Brownian motion was formally established by
\citeauthor{gatheral2018volatility}~\cite{gatheral2018volatility},
who demonstrated that realised variance increments possess a scaling
exponent $H \approx 0.07$ across a wide universe of equity indices
and maturities.  This finding has profound implications for derivative
pricing: the implied volatility surface generated by rough models
reproduces the characteristic \emph{at-the-money skew} of short-dated
options without resorting to jumps or stochastic interest rates.

Classical stochastic volatility models—Heston \cite{heston1993closed},
SABR \cite{hagan2002managing}, and their extensions—assume that the
variance process is a semi-martingale.  The empirical Hurst exponent
$H \approx 0.07 < 0.5$ violates this assumption, necessitating a
framework based on \emph{fractional} stochastic calculus.  The Rough
Heston model of \citeauthor{el2019roughening}~\cite{el2019roughening}
and the rough Bergomi model of
\citeauthor{bayer2016pricing}~\cite{bayer2016pricing} are the two
canonical representatives of this class.

The present paper makes the following contributions:

\begin{enumerate}[leftmargin=*, label=\textbf{C\arabic*.}]
  \item A production-grade implementation of the \emph{Hybrid Scheme}
        \cite{bennedsen2017hybrid} for discrete simulation of the
        stochastic Volterra equation, accelerated via Numba JIT
        compilation to achieve sub-millisecond per-path throughput.

  \item A fully vectorised Gaussian HMM calibrated by the
        Baum--Welch algorithm with multiple random restarts, whose
        forward-filtered state probabilities gate position sizing in
        real time.

  \item An Interactive Brokers connectivity layer implementing
        passive limit-order execution with a \emph{chase} algorithm
        that adjusts the limit price at configurable intervals until
        filled or a market order is issued.

  \item A scientific validation suite comprising a Monte Carlo
        convergence test and a regime stability analysis against
        historical market crises.
\end{enumerate}

The remainder of the paper is organised as follows.
Section~\ref{sec:math} establishes the mathematical framework.
Section~\ref{sec:method} details the algorithmic methodology.
Section~\ref{sec:results} presents the numerical results.
Section~\ref{sec:discussion} discusses limitations and extensions.
Section~\ref{sec:conclusion} concludes.

% ═══════════════════════════════════════════════════════════════════
\section{Mathematical Framework}\label{sec:math}

\subsection{Fractional Brownian Motion}

\begin{definition}[Fractional Brownian Motion]
A fractional Brownian motion (fBM) $B^H = \{B^H_t, t \geq 0\}$ with
Hurst exponent $H \in (0,1)$ is the unique (up to a constant) centred
Gaussian process with continuous sample paths satisfying the covariance
structure
\begin{equation}\label{eq:fbm_cov}
  \E\bigl[B^H_t\, B^H_s\bigr]
  = \frac{1}{2}\Bigl(t^{2H} + s^{2H} - \abs{t-s}^{2H}\Bigr),
  \quad t,s \geq 0.
\end{equation}
\end{definition}

When $H = 1/2$, \eqref{eq:fbm_cov} reduces to the covariance of
standard Brownian motion, $\min(t,s)$.  For $H < 1/2$, successive
increments exhibit \emph{negative} correlation—a property known as
\emph{anti-persistence} or roughness.  The increments of fBM (the
\emph{fractional Gaussian noise}) have covariance
\begin{equation}\label{eq:fgn_cov}
  \gamma(k) = \frac{1}{2}\Bigl(\abs{k+1}^{2H} - 2\abs{k}^{2H}
              + \abs{k-1}^{2H}\Bigr), \quad k \in \mathbb{Z}.
\end{equation}
For $H < 1/2$, $\gamma(k) < 0$ for all $k \geq 1$, confirming that
the process has memory of the opposite sign from that in long-memory
($H > 1/2$) processes.

\subsection{The Stochastic Volterra Equation}

The spot variance process $\{v_t\}$ is modelled as the solution to
the stochastic Volterra integral equation
\begin{equation}\label{eq:volterra}
  v_t = v_0 + \frac{1}{\Gamma\!\left(H + \tfrac{1}{2}\right)}
        \int_0^t (t-s)^{H - 1/2}\, \lambda(v_s)\, \dW_s,
\end{equation}
where $W$ is a standard Brownian motion on a filtered probability
space $(\Omega, \F, \{\F_t\}, \Prob)$, $\G$ denotes the Euler Gamma
function, and $\lambda : \R_+ \to \R_+$ is a Lipschitz diffusion
coefficient.  In the rough Bergomi specification,
$\lambda(v) = \nu\,\sqrt{v}$ for a constant $\nu > 0$ (the
\emph{vol-of-vol}).  The kernel $K(t,s) = (t-s)^{H-1/2}/\G(H+1/2)$
is integrable and defines a fractional integral operator; it is
precisely the Mandelbrot--Van Ness kernel connecting fBM to the
Wiener integral.

\begin{remark}[Singularity of the kernel]\label{rem:singular}
For $H < 1/2$, the exponent $H - 1/2 < 0$, so the kernel
$K(t,s) \to +\infty$ as $s \to t^-$.  This divergence violates the
Lipschitz continuity requirement of standard It\^{o} integration and
renders the classical Euler--Maruyama (EM) scheme inconsistent:
applying EM na\"ively assigns the full singular weight to the most
recent Brownian increment, yielding a strong approximation error of
$O(n^H)$ rather than the classical $O(n^{1/2})$
\cite{richard2021discretisation}.  For $H = 0.07$, this is
catastrophically slow—achieving the same accuracy as EM at $H = 0.5$
with $n$ steps would require $n^{0.5/0.07} \approx n^{7.1}$ steps
under the rough setting.
\end{remark}

\subsection{The Hybrid Convolution Scheme}

The \emph{Hybrid Scheme} of \citeauthor{bennedsen2017hybrid}
\cite{bennedsen2017hybrid} resolves the singularity by decomposing the
kernel into a \emph{near-field} component, handled with exact
power-law quadrature weights over $\kappa$ lags, and a
\emph{far-field} tail, approximated by a truncated exponential sum of
$J$ terms:
\begin{equation}\label{eq:hybrid}
  K(t,s) \approx \underbrace{\sum_{j=0}^{\kappa-1} w_j
                 \,\mathbf{1}_{[j\Delta t,\,(j+1)\Delta t)}(s)}_{%
                 \text{near-field (exact)}}
               + \underbrace{\sum_{l=1}^{J} c_l\,
                 e^{-\gamma_l(t-s)}}_{\text{far-field (approx.)}},
\end{equation}
where the weights $w_j = \bigl[((j+1)\Delta t)^{H+1/2} -
(j\Delta t)^{H+1/2}\bigr] / \bigl[(H+1/2)\,\G(H+1/2)\bigr]$ are
derived by exact integration of the power law over each interval, and
the exponential coefficients $(c_l, \gamma_l)$ are determined by a
least-squares fit on a geometric grid of sample points.  This
decomposition reduces the per-step computational complexity from
$O(N^2)$ (full convolution) to $O(N\kappa + NJ)$, enabling
simulation of $10{,}000$ paths in sub-second wall time on modern
hardware.

\subsection{The Gaussian HMM}

The G-HMM posits that observed log-returns $r_t = \log(S_t/S_{t-1})$
are generated by a two-state hidden Markov chain $\{s_t\}_{t=1}^T$
with state space $\mathcal{S} = \{0\,(\text{Calm}),\,
1\,(\text{Turbulent})\}$ and emission distributions
\begin{equation}\label{eq:emission}
  r_t \mid s_t = k \;\sim\; \N(\mu_k,\, \sigma_k^2), \quad k \in \{0,1\}.
\end{equation}
The complete model is parameterised by $\theta = (\bm{\pi}, A,
\bm{\mu}, \bm{\sigma})$, where $\bm{\pi}$ is the initial-state
distribution and $A = [a_{jk}]$ is the row-stochastic transition
matrix.  By design, the Turbulent state is identified as the one with
the higher conditional variance: $\sigma_1^2 > \sigma_0^2$.

\subsection{The Kelly Criterion}

Under a continuous-time approximation, the fraction $f^*$ of capital
allocated to the strategy that maximises the expected logarithmic
growth rate (the Kelly criterion) is
\begin{equation}\label{eq:kelly}
  f^* = \frac{\hat{\mu} - r_f}{\hat{\sigma}^2},
\end{equation}
where $\hat{\mu}$ and $\hat{\sigma}^2$ are rolling estimates of the
strategy's mean and variance of excess returns, and $r_f$ is the
risk-free rate.  In practice, the allocation is capped at
$f^*_{\text{cap}} = \min(f^*,\, f_{\max})$ with $f_{\max} = 0.5$
(half-Kelly) to mitigate the sensitivity of the full-Kelly criterion
to estimation error \cite{thorp2011kelly}.

% ═══════════════════════════════════════════════════════════════════
\section{Methodology}\label{sec:method}

\subsection{Simulation Algorithm}

Algorithm~\ref{alg:mc} summarises the Hybrid Scheme Monte Carlo
procedure used to price the ATM straddle.

\begin{algorithm}[h]
\caption{Hybrid Scheme Monte Carlo Straddle Pricer}\label{alg:mc}
\begin{algorithmic}[1]
  \Require $v_0, H, \nu, \rho, S_0, r, T, \kappa, J, N$
  \State Pre-compute near-field weights $\{w_j\}_{j=0}^{\kappa-1}$
         via \eqref{eq:hybrid}
  \State Fit far-field exponential coefficients $\{(c_l,\gamma_l)\}$
         by regularised least squares
  \For{$p = 1, \ldots, N$} \hfill (parallelise with \texttt{prange})
    \State Initialise $v \leftarrow v_0$, $\log S \leftarrow \log S_0$,
           circular buffer $\mathbf{b} \leftarrow \mathbf{0}_\kappa$,
           far-field state $\mathbf{x} \leftarrow \mathbf{0}_J$
    \For{$i = 0, \ldots, n-1$}
      \State Sample $(Z_1, Z_2) \sim \N(\mathbf{0}, I_2)$
      \State $\dd W_v \leftarrow \sqrt{\Delta t}(\rho Z_1
             + \sqrt{1-\rho^2}\,Z_2)$,
             $\dd W_S \leftarrow \sqrt{\Delta t}\,Z_1$
      \State $\sigma_v \leftarrow \nu\sqrt{\max(v,0)}$
      \State Near: $\eta \leftarrow \sum_{j=0}^{\min(i,\kappa)-1}
             w_j\,\mathbf{b}_{(i-j)\bmod\kappa}$
      \State Far: $x_l \leftarrow e^{-\gamma_l\Delta t}x_l
             + c_l\,\sigma_v\,\dd W_v$; $\;\eta \mathrel{+}=
             \sum_l x_l$
      \State $v \leftarrow \max(v_0 + \eta,\, 0)$
      \State $\log S \mathrel{+}= (r - v/2)\Delta t
             + \sqrt{\max(v,0)}\;\dd W_S$
      \State Push $\sigma_v\,\dd W_v$ into circular buffer
    \EndFor
    \State $\text{payoff}_p \leftarrow \abs{S_0 e^{\log S} - K}$
  \EndFor
  \State \Return $e^{-rT}\,\bar{\text{payoff}}$, $e^{-rT}\,s_{\text{payoff}}$
\end{algorithmic}
\end{algorithm}

\subsection{Baum--Welch Calibration}

The G-HMM parameters are estimated by maximising the complete-data
log-likelihood via the EM algorithm.  At the E-step, the
\emph{forward variable} $\alpha_t(k) = \Prob(r_{1:t}, s_t=k\mid\theta)$
is computed recursively:
\begin{align}
  \alpha_0(k) &= \pi_k\,\N(r_1;\mu_k,\sigma_k^2), \nonumber\\
  \alpha_t(k) &= \N(r_t;\mu_k,\sigma_k^2)
                \sum_{j} \alpha_{t-1}(j)\,a_{jk}, \label{eq:forward}
\end{align}
and the backward variable $\beta_t(k)$ is computed symmetrically.
All recursions are implemented in scaled form to prevent numerical
underflow.  At the M-step, the transition matrix and emission
parameters are updated in closed form.  Multiple random restarts with
K-means++ initialisation are employed to mitigate convergence to
local optima.

\subsection{Real-Time Regime Inference}

For live trading, the causal (forward-filtered) probability is used to
avoid look-ahead bias:
\begin{equation}\label{eq:filtering}
  P(s_t = \text{Turbulent}\mid r_{1:t}) =
  \frac{\alpha_t(\text{Turbulent})}{\alpha_t(0) + \alpha_t(1)}.
\end{equation}
The \emph{traffic-light} rule is:
\begin{itemize}
  \item $P > 0.8$: \textbf{CLOSE\_ALL} — flatten all positions immediately.
  \item $0.6 < P \leq 0.8$: \textbf{DELTA\_HEDGE} — reduce exposure
        and hedge residual delta.
  \item $P \leq 0.6$: \textbf{TRADE} — proceed with full Kelly sizing.
\end{itemize}

\subsection{Passive Execution: Chase Algorithm}

To minimise slippage on option executions, a \emph{passive-limit chase}
procedure is employed:
\begin{enumerate}
  \item Post a limit order at the midpoint
        $m = \frac{\text{bid} + \text{ask}}{2}$.
  \item If unfilled after $\tau_{\text{chase}}$ seconds, move the
        limit one tick toward the aggressive side.
  \item Repeat up to $n_{\text{chase}}$ times.
  \item If still unfilled: optionally convert to a market order.
\end{enumerate}
The tick size follows CBOE conventions: \$0.01 for options priced
below \$3.00, and \$0.05 otherwise.  Orders are placed on a dedicated
\texttt{clientId} connection so market-data and execution callbacks
never contend for the same socket.

\subsection{Net PnL and Transaction Costs}

The net daily PnL is
\begin{equation}\label{eq:net_pnl}
  \Pi_t^{\text{net}} = f_t\, r_t^{\text{strad}}
    - (c_{\text{spread}} + c_{\text{slip}})\,\abs{\Delta f_t},
\end{equation}
where $f_t$ is the Kelly fraction, $r_t^{\text{strad}}$ is the
daily straddle return, $c_{\text{spread}}$ is the bid--ask
half-spread, and $c_{\text{slip}}$ is the market-impact slippage,
both expressed as a fraction of notional.

% ═══════════════════════════════════════════════════════════════════
\section{Numerical Results and Empirical Findings}\label{sec:results}

\subsection{Roughness Visualisation}

Fig.~\ref{fig:fbm} presents simulated fBM paths for $H = 0.07$ (rough,
matching the empirical SPX Hurst exponent) and $H = 0.5$ (classical
Brownian motion).  It was observed that the rough paths exhibit
substantially more erratic local behaviour, consistent with the
anti-persistence property of $\gamma(k) < 0$ for all positive lags
(cf.~\eqref{eq:fgn_cov}).  The middle panels display rolling realised
volatility: the rough regime ($H = 0.07$) produces highly irregular,
spiky volatility clusters, whereas the standard Brownian case
maintains a relatively stable level.  The bottom row visualises the
Volterra instantaneous volatility $\sqrt{v_t}$, confirming that the
Hybrid Scheme correctly preserves the roughness of the driving kernel.
These qualitative features are in accordance with the empirical
findings of
\citeauthor{gatheral2018volatility}~\cite{gatheral2018volatility} for
the S\&P~500 index.

\begin{figure}[h]
  \centering
  \includegraphics[width=\columnwidth]{fbm_roughness.png}
  \caption{%
    \textbf{Roughness Comparison.}
    Top row: seven sample paths of fBM with $H=0.07$ (left, red tones,
    empirical SPX exponent) and $H=0.5$ (right, blue tones) over
    $[0,1]$ with $n=1{,}000$ steps, simulated via the Davies--Harte
    exact circulant embedding method.
    Middle row: the corresponding rolling realised volatility (50-step
    window); the rough process ($H=0.07$) exhibits visually apparent
    anti-persistence and clustered volatility spikes absent from the
    standard Brownian case.
    Bottom row: sample paths of the Volterra instantaneous volatility
    $\sqrt{v_t}$ generated by the discrete Hybrid Scheme; the
    rough kernel ($H=0.07$) produces substantially more erratic
    variance dynamics than the classical semi-martingale ($H=0.5$).
  }
  \label{fig:fbm}
\end{figure}

\subsection{Monte Carlo Convergence}

The convergence of the straddle price estimator is characterised in
Fig.~\ref{fig:convergence}.  Twenty-five independent Monte Carlo runs
were executed for each $N \in \{32, 64, 128, 256, 512, 1024, 2048,
4096\}$ paths.  The empirical standard error $\widehat{\mathrm{SE}}(N)
:= \mathrm{std}(\hat{P}_1, \ldots, \hat{P}_{25})$ was fitted by
ordinary least squares on the log--log scale:
\begin{equation}
  \log \widehat{\mathrm{SE}}(N) = a + b\,\log N.
\end{equation}
The fitted slope was found to be $b \approx -0.50$ with $R^2 > 0.97$,
confirming that the estimator converges at the canonical Monte Carlo
rate $O(N^{-1/2})$ and that no pathological variance inflation is
present.  A slope outside the acceptance interval
$[-0.65, -0.35]$ would be flagged as a \emph{Variance Reduction
Failure} by the automated validation suite.

\begin{figure}[h]
  \centering
  \includegraphics[width=\columnwidth]{convergence_plot.png}
  \caption{%
    \textbf{Monte Carlo Convergence.}
    Left: log--log plot of empirical standard error $\widehat{\mathrm{SE}}(N)$
    vs.\ number of paths $N$.  The fitted regression line (red dashed)
    has slope $b \approx -0.50$, indistinguishable from the theoretical
    $O(N^{-1/2})$ reference (green dotted).
    Right: absolute error $\abs{\hat{P}(N) - P_{\mathrm{BS}}}$, where
    $P_{\mathrm{BS}}$ is the Black--Scholes reference price obtained
    by setting $H = 0.5$ and $\nu \to 0$.  Both panels confirm
    asymptotically correct behaviour of the Hybrid Scheme estimator.
  }
  \label{fig:convergence}
\end{figure}

\subsection{Regime Detection}

The G-HMM was fitted to 1{,}260 synthetic daily log-returns
(five years) containing five embedded turbulent episodes generated
from $\N(-0.0015,\,0.025^2)$ against a calm baseline of
$\N(+0.0005,\,0.007^2)$.  The calibrated parameters are summarised
in Table~\ref{tab:hmm}.

\begin{table}[h]
  \centering
  \caption{Fitted G-HMM Parameters}\label{tab:hmm}
  \begin{tabular}{lcccc}
    \toprule
    State & $\hat{\mu}$ & $\hat{\sigma}$ & $\hat{\pi}$ & $\hat{a}_{kk}$ \\
    \midrule
    Calm      & $+0.00051$ & $0.00702$ & $0.803$ & $0.9994$ \\
    Turbulent & $-0.00148$ & $0.02487$ & $0.197$ & $0.9948$ \\
    \bottomrule
  \end{tabular}
\end{table}

The expected regime durations, $1/(1-\hat{a}_{kk})$, are
approximately 1{,}667 and 192 trading days for the Calm and Turbulent
states, respectively, consistent with the generating process.
The Viterbi-decoded state sequence achieved a classification
accuracy of $94.3\%$ against the true labels.

\begin{figure}[h]
  \centering
  \includegraphics[width=\columnwidth]{regime_map.png}
  \caption{%
    \textbf{Regime Map.}
    Top: daily log-returns coloured by the Viterbi-decoded state
    (blue = Calm, red = Turbulent); shaded bands indicate the true
    turbulent episodes.  Middle: forward-filtered posterior probability
    $P(s_t = \text{Turbulent} \mid r_{1:t})$ with the $0.60$ threshold
    (dashed).  Bottom: 21-day rolling annualised volatility.
    Right panel: return distributions by state with fitted Gaussian
    overlays; bottom-right: estimated transition probability matrix.
    It was observed that the posterior probability uniformly exceeded
    $0.60$ within five trading days of each embedded crisis onset.
  }
  \label{fig:regime}
\end{figure}

\subsection{Strategy Performance}

The cumulative equity curves and drawdown profiles of the two
strategies are displayed in Fig.~\ref{fig:equity}.
The regime-adjusted strategy outperformed the naive long-volatility
benchmark on all risk-adjusted metrics, as reported in
Table~\ref{tab:perf}.

\begin{table}[h]
  \centering
  \caption{Strategy Performance Summary (5-Year Simulation)}\label{tab:perf}
  \begin{tabular}{lcc}
    \toprule
    Metric & Naive & Regime-Adjusted \\
    \midrule
    Annualised Return  & $5.1\%$  & $7.3\%$  \\
    Annualised Sharpe  & $0.87$   & $1.38$   \\
    Maximum Drawdown   & $-22.4\%$& $-12.7\%$\\
    Calmar Ratio       & $0.23$   & $0.57$   \\
    \bottomrule
  \end{tabular}
\end{table}

The empirical results suggest that the HMM traffic-light mechanism
is responsible for the majority of the improvement: during turbulent
episodes, the regime-adjusted strategy reduces its position to zero,
thereby avoiding the acute drawdowns suffered by the naive strategy.
The drawdown analysis panel of Fig.~\ref{fig:equity} confirms that
the maximum drawdown of the regime-adjusted strategy occurs
exclusively outside the shaded turbulent intervals.

\begin{figure}[h]
  \centering
  \includegraphics[width=\columnwidth]{equity_curve.png}
  \caption{%
    \textbf{Strategy Performance.}
    Top: cumulative equity curves for the Naive Long-Vol strategy
    (grey) and the Regime-Adjusted strategy (green) over 1{,}260
    simulated trading days; shaded red bands denote embedded
    turbulent episodes.
    Middle: drawdown profiles; the regime-adjusted maximum drawdown
    is confined to calm periods.
    Bottom: dynamic position sizing via the Kelly criterion;
    the HMM filter sets the position to zero during turbulent episodes.
  }
  \label{fig:equity}
\end{figure}

\subsection{Transaction Cost Sensitivity}

The Sharpe ratio sensitivity to bid--ask spread $c_{\text{spread}}
\in [0, 2\%]$ and slippage $c_{\text{slip}} \in [0, 1\%]$ is
presented as a heatmap in Fig.~\ref{fig:heatmap}.  The
\emph{breakeven frontier} (contour where $\text{SR} = 0$) is clearly
delineated by the black contour line.  It was observed that the
regime-adjusted strategy maintains a positive Sharpe ratio across a
significantly wider cost region than the naive strategy, owing to its
lower portfolio turnover during turbulent regimes.  At representative
combined costs of $c_{\text{spread}} = 0.5\%$ and $c_{\text{slip}} =
0.2\%$, the regime-adjusted Sharpe exceeds that of the naive strategy
by $0.42$ Sharpe units.

\begin{figure}[h]
  \centering
  \includegraphics[width=\columnwidth]{transaction_cost_heatmap.png}
  \caption{%
    \textbf{Transaction Cost Analysis.}
    Sharpe ratio as a function of bid--ask spread ($x$-axis) and
    slippage ($y$-axis) for: Naive Long-Vol (left), Regime-Adjusted
    (centre), and the difference (right).  Red/green colouring
    indicates negative/positive Sharpe; the black contour is the
    breakeven frontier ($\text{SR} = 0$).  The regime-adjusted
    strategy retains a positive Sharpe across a substantially larger
    feasible cost region.
  }
  \label{fig:heatmap}
\end{figure}

% ═══════════════════════════════════════════════════════════════════
\section{Discussion}\label{sec:discussion}

\subsection{Stability of the Rolling Hurst Estimation}

The rolling Hurst estimator employed in the live system uses
Detrended Fluctuation Analysis (DFA) over a trailing window of
$\sim\!252$ trading days.  It was noted that the estimator exhibits
non-trivial finite-sample bias for $H < 0.1$: the DFA fluctuation
function is susceptible to trending artefacts at short timescales,
causing the estimated $\hat{H}$ to drift upward by approximately
 +0.03$ relative to the true value.  This bias is conservative for
the purposes of the pricing engine—overestimating $H$ makes the
simulated kernel \emph{less} singular—but it warrants periodic
re-calibration with out-of-sample data.

\subsection{Efficacy of the $P > 0.6$ Risk Threshold}

The choice of $0.6$ as the turbulence threshold is motivated by a
bias-variance trade-off: lower thresholds increase the frequency of
risk-off signals, reducing drawdowns at the cost of forfeited positive
carry in calm markets; higher thresholds preserve more upside but
expose the portfolio to regime transitions.  A grid search over the
synthetic dataset suggests that $0.6$ is near-optimal in terms of
Sharpe ratio, but the optimum shifts to approximately $0.55$ under
elevated transaction costs—a regime in which frequent position changes
are more penalised.  The sensitivity of this finding to the underlying
model parameters (vol-of-vol $\nu$, correlation $\rho$) and to
sample-path realisations remains a direction for further investigation.

\subsection{Limitations}

Several limitations of the present framework merit acknowledgment.
First, the G-HMM assumes Gaussian emission distributions; empirical
equity returns exhibit heavier tails, which may lead to systematic
misclassification of extreme events.  A Student-$t$ or
asymmetric Laplace emission distribution could improve robustness.
Second, the pricing model is calibrated to a fixed Hurst exponent;
in practice, $H$ varies across market regimes, and a
\emph{regime-conditional} $H$ would improve option pricing accuracy.
Third, the backtest is conducted on synthetic data; out-of-sample
validation on historical SPY options data is required before
any capital commitment.

% ═══════════════════════════════════════════════════════════════════
\section{Comprehensive Quantitative Model Taxonomy}\label{sec:taxonomy}

It is instructive to situate the Rough-Regime framework within the
broader landscape of stochastic volatility models.  The present
section provides a systematic comparative analysis of four canonical
models—Heston, SABR, Local Volatility (Dupire), and Hawkes
Jump-Diffusion—evaluated against the Volterra process that constitutes
the pricing core of the present system.

\subsection{Heston Stochastic Volatility}

The Heston (1993) model \cite{heston1993closed} specifies the joint
risk-neutral dynamics of the spot price $S_t$ and its instantaneous
variance $V_t$ via the coupled system of SDEs:
\begin{align}
  dS_t &= r\,S_t\,\dd t + \sqrt{V_t}\,S_t\,\dd W_t^S,
  \label{eq:heston_s}\\
  dV_t &= \kappa(\theta - V_t)\,\dd t + \xi\,\sqrt{V_t}\,\dd W_t^V,
  \label{eq:heston_v}
\end{align}
where $d\langle W^S, W^V\rangle_t = \rho\,\dd t$.  The parameters
$\kappa > 0$, $\theta > 0$, and $\xi > 0$ denote the mean-reversion
speed, long-run variance, and vol-of-vol, respectively.  The
\emph{Feller condition} $2\kappa\theta > \xi^2$ ensures that $V_t$
remains strictly positive almost surely, precluding degenerate
zero-volatility states.

The Heston model admits a semi-closed-form characteristic function,
enabling rapid option pricing via Fourier inversion.  Negative
correlation $\rho < 0$ induces the empirically observed
\emph{leverage effect}: a negative return in the spot is accompanied
by an increase in variance, generating a downward IV skew.  However,
since $V_t$ is a continuous semi-martingale (Hurst exponent $H = 0.5$),
the model cannot reproduce the ultra-steep short-dated ATM skew
observed empirically for $H \approx 0.07$.

\begin{figure}[h]
  \centering
  \includegraphics[width=\columnwidth]{heston_surface.png}
  \caption{%
    \textbf{Heston Implied Volatility Surface.}
    Left: the three-dimensional Monte Carlo implied volatility surface
    for the Heston model with parameters $\kappa=2.0$, $\theta=0.04$,
    $\xi=0.3$, and $\rho=-0.7$.  The negative correlation parameter
    induces a pronounced downward skew (leverage effect) that is more
    pronounced at shorter maturities.  Right: fifty sample paths of
    the instantaneous volatility $\sqrt{V_t}$ mean-reverting toward
    $\sqrt{\theta}=20\%$ (dashed).
  }
  \label{fig:heston}
\end{figure}

\subsection{SABR Stochastic Volatility}

The SABR model \cite{hagan2002managing} was developed to address the
shortcomings of earlier stochastic volatility models in the interest-rate
derivatives context.  It specifies the dynamics of a forward rate $F_t$
and its volatility $\sigma_t$ as:
\begin{align}
  dF_t      &= \sigma_t\,F_t^\beta\,\dd W_t^1,
  \label{eq:sabr_f}\\
  d\sigma_t &= \alpha\,\sigma_t\,\dd W_t^2,
  \label{eq:sabr_sigma}
\end{align}
where $d\langle W^1, W^2\rangle_t = \rho\,\dd t$ and
$\beta \in [0, 1]$ interpolates between the normal ($\beta=0$) and
log-normal ($\beta=1$) backbones.  The parameter $\alpha > 0$
controls the vol-of-vol, and $\nu := \alpha$ in the Hagan notation.

The principal utility of SABR derives from its \emph{analytic
approximation} for the Black implied volatility
$\sigma_{\text{imp}}(F, K, T)$, which takes the form \cite{hagan2002managing}:
\begin{equation}\label{eq:sabr_iv}
  \sigma_{\text{imp}} \approx
  \frac{\alpha}{(FK)^{(1-\beta)/2}}\cdot
  \frac{z}{\chi(z)}\cdot
  \Bigl[1 + \Bigl(\tfrac{(1-\beta)^2}{24}\tfrac{\alpha^2}{(FK)^{1-\beta}}
  + \tfrac{\rho\beta\nu\alpha}{4(FK)^{(1-\beta)/2}}
  + \tfrac{2-3\rho^2}{24}\nu^2\Bigr)T\Bigr],
\end{equation}
where $z = \frac{\nu}{\alpha}(FK)^{(1-\beta)/2}\ln(F/K)$ and
$\chi(z)=\ln\!\frac{\sqrt{1-2\rho z+z^2}+z-\rho}{1-\rho}$.
This approximation enables rapid calibration to the entire swaption
or cap volatility cube and constitutes the de facto standard in
fixed-income options markets.

\begin{figure}[h]
  \centering
  \includegraphics[width=\columnwidth]{sabr_smile.png}
  \caption{%
    \textbf{SABR Implied Volatility Smile.}
    Left: implied volatility as a function of strike for four
    SABR parameter configurations at one-year maturity with forward
    rate $F=3\%$, illustrating the symmetric smile ($\rho=0$),
    downward skew ($\rho=-0.4$), high-vol-vol curvature ($\nu=0.8$),
    and the log-normal backbone ($\beta=1$).
    Right: ATM implied volatility term structure for varying vol-of-vol
    levels, demonstrating the model's ability to replicate the
    observed flattening of the vol term structure at long maturities.
  }
  \label{fig:sabr}
\end{figure}

\subsection{Local Volatility and the Dupire Equation}

Dupire (1994) \cite{dupire1994pricing} and Derman--Kani (1994)
\cite{derman1994riding} independently established that any
arbitrage-free call price surface $C(T, K)$ is consistent with a
\emph{unique} Markovian diffusion of the form
$dS_t = r\,S_t\,\dd t + \sigma_{\text{loc}}(t,S_t)\,S_t\,\dd W_t$,
where the deterministic local volatility function satisfies the
\emph{Dupire equation}:
\begin{equation}\label{eq:dupire}
  \sigma_{\text{loc}}^2(T,K) =
  \frac{\partial_T C + r K\,\partial_K C}
       {\frac{1}{2}K^2\,\partial_{KK}C}.
\end{equation}
Equation~\eqref{eq:dupire} is an \emph{exact} inversion: by
construction, the local volatility model reproduces every observable
vanilla option price without error.  This property renders local
volatility the canonical model for the valuation of path-dependent
exotic instruments, including barrier options, Asian options, and
lookback contracts, for which static replication arguments are
unavailable.

The practical limitation of local volatility is that the
\emph{forward volatility smile}—the implied volatility surface
conditional on reaching a future state—is necessarily flat, in
contradistinction to the empirical term structure of the smile.  This
deficiency motivates the stochastic local volatility (SLV) extension,
in which the Dupire surface is used to calibrate a mixing fraction
between the local and Heston components.

\begin{figure}[h]
  \centering
  \includegraphics[width=\columnwidth]{dupire_local_vol.png}
  \caption{%
    \textbf{Dupire Local Volatility Calibration.}
    Left: the synthetic market implied volatility surface exhibiting
    a pronounced downward skew and flattening term structure consistent
    with the SPX volatility surface.
    Right: the corresponding Dupire local volatility surface calibrated
    via numerical differentiation of the call price surface.  The
    local vol surface is more steeply sloped and exhibits stronger
    short-maturity skew than the implied surface, consistent with the
    theoretical relationship $\sigma_{\text{loc}} \approx
    2\,\sigma_{\text{imp}}\,(\cdot)$ in the small-vol-of-vol limit.
  }
  \label{fig:dupire}
\end{figure}

\subsection{Hawkes Jump-Diffusion and Market Contagion}

The self-exciting Hawkes process \cite{hawkes1971spectra} provides a
mathematically rigorous mechanism for the \emph{contagion} and
\emph{clustering} of extreme market events.  The conditional intensity
of a univariate Hawkes process is defined as:
\begin{equation}\label{eq:hawkes_intensity}
  \lambda(t) = \mu + \sum_{t_i < t} \phi(t - t_i),
  \qquad \phi(s) = \alpha_H\,e^{-\beta_H s},
\end{equation}
where $\mu > 0$ is the baseline (exogenous) intensity, $\alpha_H > 0$
is the excitation coefficient (the instantaneous jump in intensity
caused by each event), and $\beta_H > 0$ is the decay rate.  The
process is stationary if and only if $\alpha_H / \beta_H < 1$,
which bounds the expected number of offspring per event.

The Hawkes Jump-Diffusion price process is then:
\begin{equation}\label{eq:hawkes_jd}
  dS_t = \mu_S\,S_t\,\dd t + \sigma_S\,S_t\,\dd W_t
         + S_{t^-}\bigl(e^J - 1\bigr)\,\dd N_t,
\end{equation}
where $J \sim \N(\mu_J, \sigma_J^2)$ and $N_t$ is the Hawkes counting
process with intensity \eqref{eq:hawkes_intensity}.  Negative
$\mu_J < 0$ models the stylised fact that crash events are
predominantly downward jumps.

\begin{figure}[h]
  \centering
  \includegraphics[width=\columnwidth]{hawkes_jump_diffusion.png}
  \caption{%
    \textbf{Hawkes Jump-Diffusion Simulation.}
    Top row: six price paths under the Hawkes jump-diffusion model;
    vertical orange dashes on the representative path (red) indicate
    jump times, which are visibly clustered.
    Middle row, left: the conditional intensity $\lambda(t)$, exhibiting
    self-exciting spikes following each event.  Middle row, right:
    the clustering of jump times across six paths, demonstrating the
    contagion dynamics absent from Poisson jump-diffusion models.
    Bottom row, left: sensitivity of intensity to the excitation
    parameter $\alpha_H$.  Bottom row, right: the return distribution
    is leptokurtic relative to the Gaussian fit, a direct consequence
    of the clustered jump mechanism.
  }
  \label{fig:hawkes}
\end{figure}

\subsection{HMM--Hawkes Integration and Regime-Conditioned Contagion}

A key theoretical contribution of the present framework is the
identification of a natural hierarchical coupling between the
macro-level G-HMM regime filter and the micro-level Hawkes jump
process.  The G-HMM operates on the daily return series to infer
the posterior probability of residing in the Turbulent state,
$P(s_t = \text{Turbulent} \mid r_{1:t})$, as described in
Section~\ref{sec:method}.  This posterior probability may be
interpreted as a \emph{regime-conditional prior} on the baseline
intensity $\mu$ of the Hawkes process: when $P > 0.6$, the
exogenous jump rate is elevated, increasing both the expected cluster
size and the tail-risk of the position.

\begin{figure}[h]
  \centering
  \includegraphics[width=\columnwidth]{hmm_hawkes_integration.png}
  \caption{%
    \textbf{HMM--Hawkes Integration.}
    Left: the forward-filtered posterior probability of the Turbulent
    state under the G-HMM, with the $0.6$ (dashed red) and $0.8$
    (dotted dark-red) risk thresholds indicated.  Shaded regions
    denote the three embedded turbulent episodes.
    Right: the Hawkes conditional intensity $\lambda(t)$ during each
    turbulent episode, illustrating the self-exciting cascade dynamics
    that emerge within periods of elevated macro-regime risk.  The
    HMM serves as an \emph{early-warning indicator}, while the Hawkes
    process quantifies the within-crisis contagion microstructure.
  }
  \label{fig:hmm_hawkes}
\end{figure}

It is thus proposed that a fully integrated system would operate in
two layers: (i) the G-HMM classifies the prevailing macro regime and
modulates position size via the traffic-light rule, and (ii) a
regime-conditional Hawkes process, calibrated separately in each HMM
state, provides a real-time jump intensity estimate used to dynamically
adjust the tail-risk buffer required by the risk management layer.
This two-tier architecture is consistent with the empirical observation
that volatility clustering occurs at multiple timescales simultaneously
\cite{gatheral2018volatility}.

\subsection{Decision Matrix: When to Use Which Model}

Table~\ref{tab:taxonomy} provides a systematic comparison of the five
models surveyed in this work, organised by the principal market
application and the structural feature that motivates their use.

\begin{table*}[t]
  \centering
  \caption{Quantitative Model Taxonomy: Structural Properties and Primary Use Cases}
  \label{tab:taxonomy}
  \renewcommand{\arraystretch}{1.25}
  \begin{tabular}{p{2.2cm} p{2.5cm} p{2.0cm} p{2.8cm} p{3.5cm}}
    \toprule
    \textbf{Model} & \textbf{Variance Process} & \textbf{Hurst $H$} &
    \textbf{Primary Use Case} & \textbf{Key Limitation} \\
    \midrule
    Rough Volterra
      & Non-Markovian Volterra; $\lambda(v)=\nu\sqrt{v}$
      & $\approx 0.07$
      & Short-dated equity index IV skew; vol-of-vol dynamics
      & Path-simulation cost; no analytic price \\
    Heston
      & CIR mean-reverting; $dV=\kappa(\theta-V)\,dt+\xi\sqrt{V}\,dW$
      & $= 0.5$
      & Equity structured products; long-dated vol surfaces
      & Insufficient short-term skew; semi-martingale \\
    SABR
      & Log-normal; no mean reversion
      & $= 0.5$
      & IR swaptions; FX smile; cap/floor cubes
      & No mean reversion; degenerates at long $T$ \\
    Local Vol (Dupire)
      & Deterministic $\sigma(t,S)$; calibrated
      & N/A
      & Barrier options; Asian options; path-dependent exotics
      & Flat forward smile; misspecifies dynamics \\
    Hawkes JD
      & Jump intensity self-excitation; no diffusive vol
      & N/A
      & Credit events; contagion; crisis-clustering
      & No continuous volatility dynamics \\
    \bottomrule
  \end{tabular}
\end{table*}

The models are ordered by their degree of departure from the
classical semi-martingale assumption.  The Rough Volterra process
occupies a unique position: it is the only model in Table~\ref{tab:taxonomy}
that simultaneously (i) matches the empirical Hurst exponent
$H \approx 0.07$, (ii) generates the observed power-law term
structure of the ATM skew, and (iii) is compatible with a
multi-regime architecture via the G-HMM traffic-light gating
mechanism described in Section~\ref{sec:method}.  The remaining
models serve specialised roles that are either complementary
(Hawkes for crisis quantification) or appropriate for asset classes
with distinct empirical IV surface characteristics (SABR for
interest rates, Local Vol for exotics desks).

% ═══════════════════════════════════════════════════════════════════
\section{Conclusion}\label{sec:conclusion}

It has been demonstrated that the integration of a rough-volatility
Volterra pricing engine with a Gaussian Hidden Markov Model regime
filter yields a statistically superior strategy to a naive
long-volatility benchmark.  The Hybrid Scheme simulation algorithm
achieves $O(N^{-1/2})$ Monte Carlo convergence as verified
experimentally, and the HMM traffic-light mechanism reduces the
maximum drawdown by more than $40\%$ relative to constant exposure
at representative transaction cost levels.  The event-driven
Interactive Brokers connectivity layer implements passive limit-order
execution with automatic reconnection, rate limiting, and batched HDF5
tick storage, constituting a production-ready infrastructure for
volatility arbitrage strategies.

Future work will address: (i) extension to a three-state HMM to
distinguish \emph{trending} from \emph{calm} regimes; (ii)
incorporation of the VIX term structure as an additional observation
variable; (iii) empirical calibration to intraday SPY option data
using the Polygon.io API; and (iv) implementation of the
two-tier HMM--Hawkes contagion architecture introduced in
Section~\ref{sec:taxonomy}, in which regime-conditional Hawkes
intensities provide dynamic tail-risk buffers for the Kelly
position-sizing layer.

% ═══════════════════════════════════════════════════════════════════
\section*{Acknowledgement}

The authors gratefully acknowledge the numerical infrastructure
provided by the Numba~\cite{lam2015numba} and PyTables projects,
which were essential for achieving the throughput targets described
in Section~\ref{sec:method}.

% ═══════════════════════════════════════════════════════════════════
\begin{thebibliography}{99}

\bibitem{gatheral2018volatility}
J.~Gatheral, T.~Jaisson, and M.~Rosenbaum,
``Volatility is rough,''
\textit{Quantitative Finance}, vol.~18, no.~6, pp.~933--949, 2018.

\bibitem{bennedsen2017hybrid}
M.~Bennedsen, A.~Lunde, and M.~S.~Pakkanen,
``Hybrid scheme for Brownian semistationary processes,''
\textit{Finance \& Stochastics}, vol.~21, no.~4, pp.~931--965, 2017.

\bibitem{bayer2016pricing}
C.~Bayer, P.~Friz, and J.~Gatheral,
``Pricing under rough volatility,''
\textit{Quantitative Finance}, vol.~16, no.~6, pp.~887--904, 2016.

\bibitem{el2019roughening}
O.~El~Euch and M.~Rosenbaum,
``The characteristic function of rough Heston models,''
\textit{Mathematical Finance}, vol.~29, no.~1, pp.~3--38, 2019.

\bibitem{heston1993closed}
S.~L.~Heston,
``A closed-form solution for options with stochastic volatility with
  applications to bond and currency options,''
\textit{The Review of Financial Studies}, vol.~6, no.~2, pp.~327--343,
1993.

\bibitem{hagan2002managing}
P.~S.~Hagan, D.~Kumar, A.~S.~Lesniewski, and D.~E.~Woodward,
``Managing smile risk,''
\textit{Wilmott Magazine}, vol.~1, pp.~84--108, 2002.

\bibitem{richard2021discretisation}
A.~Richard, X.~Tan, and F.~Yang,
``Discrete-time simulation of Stochastic Volterra equations,''
\textit{Stochastic Processes and their Applications},
vol.~141, pp.~109--138, 2021.

\bibitem{rabiner1989tutorial}
L.~R.~Rabiner,
``A tutorial on hidden Markov models and selected applications in speech
  recognition,''
\textit{Proceedings of the IEEE}, vol.~77, no.~2, pp.~257--286, 1989.

\bibitem{hamilton1989new}
J.~D.~Hamilton,
``A new approach to the economic analysis of nonstationary time series
  and the business cycle,''
\textit{Econometrica}, vol.~57, no.~2, pp.~357--384, 1989.

\bibitem{thorp2011kelly}
E.~O.~Thorp,
``The Kelly criterion in blackjack, sports betting, and the stock
  market,''
in \textit{The Kelly Capital Growth Investment Criterion},
World Scientific, 2011, pp.~789--832.

\bibitem{lam2015numba}
S.~K.~Lam, A.~Pitrou, and S.~Seibert,
``Numba: A LLVM-based Python JIT compiler,''
in \textit{Proc.\ 2nd Workshop on the LLVM Compiler Infrastructure
in HPC}, 2015, pp.~1--6.

\bibitem{dupire1994pricing}
B.~Dupire,
``Pricing with a smile,''
\textit{Risk}, vol.~7, no.~1, pp.~18--20, 1994.

\bibitem{derman1994riding}
E.~Derman and I.~Kani,
``Riding on a smile,''
\textit{Risk}, vol.~7, no.~2, pp.~32--39, 1994.

\bibitem{hawkes1971spectra}
A.~G.~Hawkes,
``Spectra of some self-exciting and mutually exciting point processes,''
\textit{Biometrika}, vol.~58, no.~1, pp.~83--90, 1971.

\bibitem{aitsahalia2015hawkes}
Y.~A{\"i}t-Sahalia, J.~Cacho-Diaz, and R.~J.~Laeven,
``Modeling financial contagion using mutually exciting jump processes,''
\textit{Journal of Financial Economics}, vol.~117, no.~3, pp.~585--606, 2015.

\end{thebibliography}

\end{document}

